% =====================================================================
%  Summary of the final report
% =====================================================================
\section*{Summary}
\addcontentsline{toc}{section}{Summary}

This is the final report of Assignment~3 (milestone M3, specification section~6.13: \emph{improved fusion and
ablations complete}). We classify recipe web pages of UPMC Food-101 into 101 dishes from one image and the page
title. The dataset proposal (M1) and the M2 draft were approved by the instructor.

\paragraph{What was built.}
\begin{itemize}
  \item \textbf{Two unimodal baselines:} UA, an image model on ResNet-50, and UB, a text model on DistilBERT.
  \item \textbf{Three fusions:}
        \begin{itemize}
          \item F0, a late fusion of the two softmax outputs;
          \item F1, a concat fusion trained end to end;
          \item F2, the improved fusion. It is a cross-modal transformer over image-region and title tokens, with
                auxiliary unimodal losses and modality dropout.
        \end{itemize}
  \item \textbf{Two text references:} a keyword rule (K0) and TF-IDF with logistic regression (T0).
  \item \textbf{Three ablations of F2:} the auxiliary losses (E1), class names masked in every title (E2) and the
        modality dropout (E4). F2 and its ablations are also tested with each modality removed from part of the
        test set (E-mod).
\end{itemize}
All trained models were run with three seeds. Every decision uses the rule registered in the proposal.

\paragraph{Main findings.}
\begin{enumerate}
  \item \textbf{Fusion helps a lot.} Every fusion is $6$--$7$\,pp of macro-F1 above the best unimodal model
        (UB, $86.8\%$).
  \item \textbf{F2 beats F1 but not F0.} F2 reaches $93.53\pm0.07\%$ macro-F1. It is significantly better than
        F1 ($+0.73$\,pp) and as good as F0 ($-0.33$\,pp, not significant). F2 is also the best-calibrated model.
  \item \textbf{Class names drive the result.} Masking them drops the text-only model from $86.8$ to $43.7\%$.
        On masked titles, F2 is the best fusion: $+2.5$\,pp over F0 and $+0.9$\,pp over F1, both significant.
  \item \textbf{The auxiliary losses matter, the modality dropout only helps with a missing title.}
        \begin{itemize}
          \item Removing the auxiliary losses costs $0.46$\,pp (significant). It also weakens the image branch:
                accuracy without the title falls from $65.3$ to $58.6\%$.
          \item Removing the modality dropout does not change the full-test score. It costs $4.7$\,pp when the
                title is removed.
        \end{itemize}
  \item \textbf{The conflict groups explain the remaining errors.} A manual review of 100 conflict items shows
        two main causes. When only the image model is right, the title is mostly generic ($84\%$); fusion fixes
        almost all of these. When only the text model is right, the image mostly looks like another dish ($68\%$).
        The hard case is a title that names another dish: every fusion gets only about half of these right.
\end{enumerate}

\paragraph{Traceability.}
\begin{itemize}
  \item \textbf{Split:} every number comes from the committed split \texttt{data/splits/split.json}.
  \item \textbf{Code:}
        \begin{itemize}
          \item UA, UB and F1 (M2): commit \texttt{cfc509a};
          \item F2, the ablations and T0 (M3): commit \texttt{b51d3fb};
          \item analyses: commit \texttt{850b6bb} and later.
        \end{itemize}
        Each run records its commit in its \texttt{env.json} file, always with a clean working tree.
  \item \textbf{Checkpoints:} the selected checkpoints, with SHA-256 checksums, are linked in the Required Links
        section.
\end{itemize}
